\section{Implications for Interstate 2.0}

\subsection{Investment Sequencing: T1 First, T2 Relief Second}

The bottleneck cost analysis supports the tier-based investment sequencing proposed in \citet{ROUTE_E2}:

\textbf{Phase 1 priority 1 (Donner Pass)}: The \$4B freight tunnel investment has a 2.5-year payback at \$1.6B/year in avoided weather closure costs. No other I2.0 investment has this NPV profile. It should be the first infrastructure construction project in the program.

\textbf{Phase 1 priority 2 (T1 managed freight lanes, highest bottleneck density first)}: I-95 (9 ATRI locations, \$4.9B/yr combined), I-75 Atlanta (\$1.3B/yr), I-80 Bay Area (\$428M/yr). Managed freight lanes on these segments reduce congestion costs while also providing the dedicated freight capacity that enables reliable transit times.

\textbf{Phase 2 (T2 relief to reduce T1 load)}: I-285 Atlanta (\$3.0B/yr) requires T2 investment because the I-285 congestion is the primary cause of truck diversions onto I-75 and I-85 through Atlanta. Reducing I-285 congestion is equivalent to managed lane investment on the T1 corridors that serve Atlanta, at lower per-dollar cost.

\subsection{Atlanta: The Structural Exception}

The Atlanta corridor cluster presents a structural exception to the T2 deferral rule. I-285 cannot be dismissed as a "locally stressed T2 connector" when it accounts for \$3.0B in annual freight costs --- more than I-90's entire bottleneck inventory. The practical investment case is:

\begin{enumerate}
\item Upgrade I-285 to T1 designation (based on functional role, not betweenness centrality)
\item Build managed freight lanes on I-285 simultaneously with I-75 Atlanta
\item Build diamond connector roads between I-285 and I-75/I-85 to distribute the T1/T1 transfer load
\end{enumerate}

This is the argument for I-285 as a T1 candidate not identified in A.1's centrality-adjusted classification --- the bottleneck data provides the economic justification that pure network topology does not.

\subsection{Managed Lane PTI Targets by Bottleneck Severity}

The tier standards specification (T1: PTI ≤ 1.15) is a single target across all T1 corridors. The bottleneck data suggests differentiated targets are warranted:

\begin{itemize}
\item I-95 Northeast Corridor (current PTI 1.8--2.2): PTI ≤ 1.15 would eliminate \$4.9B/yr in bottleneck costs --- the highest absolute gain from reaching the standard
\item I-75 Atlanta (current PTI 1.84): PTI ≤ 1.15 would eliminate \$1.3B/yr
\item I-40 (current PTI 0.84): already at standard --- no managed lane investment needed
\end{itemize}

Corridors not at standard but close (PTI 1.15--1.3) can achieve the target with one managed lane pair; corridors far above standard (PTI > 1.5) require two to three pairs plus T2 relief investment.
